\cvsection{Open Source}

\begin{cventries}

	%---------------------------------------------------------

	\cventry
	{Maintainer}
	{\href{https://github.com/llp42/re-mach}{re-mach}}
	{Remote}
	{2026 -- Present}
	{
		\begin{cvitems}
			\item Maintain \textbf{re-mach}, a Mach microkernel reimplemented from scratch in \textbf{Rust}, targeting compatibility with \textbf{GNU Mach}.
			\item Implemented the core kernel --- \textbf{IPC, threads, tasks, virtual memory, and scheduling} --- now passing the basic \textbf{GNU Mach} test suite under \textbf{Multiboot2}.
		\end{cvitems}
	}

	\cventry
	{Contributor}
	{\href{https://github.com/rust-osdev/x86_64}{rust-osdev/x86\_64}}
	{Remote}
	{2026 -- Present}
	{
		\begin{cvitems}
			\item Contribute to \textbf{rust-osdev/x86\_64}, the low-level hardware abstraction crate used across the Rust OS-development ecosystem.
			\item Contributions span CPU instruction support, \textbf{RFLAGS} correctness at the compiler boundary, and \textbf{atomic page table types}.
		\end{cvitems}
	}

	\cventry
	{Contributor}
	{GNU Mach (GNU Hurd)}
	{Remote}
	{2005 -- Present}
	{
		\begin{cvitems}
			\item Contributing to the \textbf{GNU Mach} microkernel since 2005 --- mostly bug fixes and code cleanliness improvements across the kernel.
			\item Deep familiarity with Mach's IPC, VM, and task internals grounds the compatibility surface targeted by \textbf{re-mach}.
		\end{cvitems}
	}

	%---------------------------------------------------------
\end{cventries}
